\chapter{結言}

\section{結論}
本研究は、日本企業の有価証券報告書「事業等のリスク」欄に着目し、前年からの内容更新（新規性）を埋め込みベースの意味的類似度により定量化した上で、提出日周辺の市場反応との関係を検証した。主な結果は以下のとおりである。

第一に、リスク開示の新規性は、価格反応（CAR）よりも投資家の注目（異常出来高）と強く結びついた。すなわち、新規性の高いリスク開示ほど提出日周辺で取引が活発化する傾向が確認された。第二に、不確実性指標およびCARについては、同一の基本仕様では統計的に明確な関係が確認されず、リスク開示更新の効果が平均的な価格調整としては限定的である可能性が示唆された。第三に、カテゴリ別では、災害/BCP、規制/コンプライアンス、人材等で注目度との関係が比較的強く、リスク種類により反応が異なる可能性が示された。

\section{貢献の整理}
本研究の貢献は次の3点に整理できる。第一に、開示研究において「量」ではなく「更新（新規性）」に焦点を当て、日本の有価証券報告書リスク欄における年度間の意味的変化を扱った点である。第二に、チャンク分割と埋め込みベース類似度を用いて、言い換えや表現揺れを一定程度吸収しつつ新規性を測定する枠組みを提示した点である。第三に、市場反応を価格反応だけでなく注目度・不確実性の複数指標で検証し、リスク開示の新規性がどの側面に現れやすいかを識別した点である。

\section{今後の展望}
今後の課題として、第一に識別の強化が挙げられる。外生的ショックや制度変更、あるいは開示規制の改定等を利用し、新規性の変動がより外生的となる状況で検証することが望ましい。第二に、提出日以外の情報公開タイミングを含め、リスク情報が市場に取り込まれるプロセスをより精緻に捉える必要がある。第三に、新規性を一つのスコアに集約するだけでなく、追加・削除・具体化・強調といった更新タイプの分類により、どのタイプの更新が注目や不確実性に結びつくのかを検証することが重要である。第四に、投資家属性、ニュース環境、業種特性等による反応の異質性を体系的に分析し、外的妥当性を高めることが課題である。


